\textbf{Seats}
\begin{enumerate}
    \item 15 mins: Intro slides \par
    \textbf{Motivation}
    \item 20 mins: Finish Example 1.1.1 \par
    \textbf{Check piazza}
    \item 30 mins: Finish Example 1.1.2 \par
    \textbf{Check piazza}
    \item 38 mins: Finish Example 1.1.3 \par
    \textbf{Check piazza}
    \item 46 mins: Finish Example 1.1.4 \par
    \textbf{Check piazza}
    \item If needed: skip 1.1.5 until next time \par
    \textbf{Seats}
\end{enumerate}

\vspace{10mm}
\textbf{Example 1.1.2}

\begin{align*}
x+1 &= -\frac{1}{2}x+4\\
\frac{3}{2}x &= 3\\
x &= 2\\
y &= 3
\end{align*}

\vspace{10mm}
\textbf{Example 1.1.3} \par
Complete the square
\begin{align*}
(x+2)^2 - 4 &= x^2 + 4x\\
(y-3)^2 - 9 &= y^2 - 6y
\end{align*}
Substitute in
\begin{align*}
(x+2)^2 - 4 + (y-3)^2 - 9 &= 3\\
(x+2)^2 + (y-3)^2 &= 16
\end{align*}
Therefore $a=-2$, $b=3$, $R=4$.

\vspace{10mm}
\textbf{Example 1.1.4}\par
\begin{align*}
    (x-1)^2 + (x+1-2)^2 &= 8\\
    \therefore 2(x-1)^2 &= 8\\
    \therefore x-1 &= \pm 2\\
    \therefore x &= \{-1,\; 3\}
\end{align*}
So the crossing points are
\begin{align*}
(-1,0) \quad (3,4)
\end{align*}

\vspace{10mm}
\textbf{Example 1.1.5}
Gradient of chord
\begin{align*}
    m = \frac{2-6}{5-1} = -1
\end{align*}
So gradient of line through centre has equation
\begin{align*}
    y = x + c
\end{align*}
Centre of the chord is at 
\begin{align*}
(A+B)/2 = (3,4)
\end{align*}
sub into above
\begin{align*}
    4 &= 3 + c\\
    c &= 1
\end{align*}
So find where $y=x+1$ and $y=2x-3$ cross
\begin{align*}
    x+1 &= 2x-3\\
    x &= 4
\end{align*}
Centre is (4,5).